\section{Conclusion}
	The threat model developed for ViperVault demonstrates a strong alignment with established security and privacy standards, including OWASP, NIST, GDPR and the NIS2 directive.
	At its current stage of development, where the application operates entirely in local-only mode, the most pressing risks are those related to the end-user device itself.
	Examples include the presence of local malware, clipboard hijacking or physical device theft.
	These are threats that can be effectively mitigated with secure enclave usage, automatic vault locking and secure clipboard handling.
	Looking ahead, the planned evolution of ViperVault towards a synchronization server and browser extensions introduces a more complex threat landscape.
	In these future scenarios, the attack surface broadens and more sophisticated adversaries may come into play, such as supply-chain compromises, advanced persistent threats or side-channel attacks.
	Addressing these challenges from the design phase is essential in order to ensure that the transition to a multi-platform, cloud-enabled ecosystem does not erode the strong guarantees already present in the current mobile app.
	The proposed improvements --- such as the introduction of a duress or panic mode, the design of a secure and non-recoverable master password reset mechanism and anti-forensic measures like memory wiping and overlay detection --- further strengthen the defensive posture of the project.
	Combined with ViperVault’s commitment to security- and privacy-by-design principles, these measures provide a clear roadmap for maintaining both user trust and regulatory compliance as the application grows in scope and functionality.
	In conclusion, ViperVault’s threat model not only addresses the risks of the present but also anticipates those of the future, laying down a foundation for a robust, resilient and trustworthy password management solution.